% Generated draft: requires mathematical and source review.
\documentclass[11pt]{article}
\usepackage[margin=25mm]{geometry}
\usepackage{fontspec}
\setmainfont{Latin Modern Roman}
\usepackage{amsmath,amssymb,mathtools,unicode-math}
\setmathfont{Latin Modern Math}
\usepackage{enumitem,needspace,xurl,hyperref}
\setlength{\parindent}{0pt}
\setlength{\parskip}{5pt}
\setlength{\emergencystretch}{3em}
\newcommand{\N}{\mathbb{N}}
\newcommand{\Z}{\mathbb{Z}}
\newcommand{\Q}{\mathbb{Q}}
\newcommand{\R}{\mathbb{R}}
\newcommand{\C}{\mathbb{C}}
\newcommand{\F}{\mathbb{F}}
\newcommand{\A}{\mathbb{A}}
\newcommand{\PP}{\mathbb{P}}
\newcommand{\E}{\mathbb{E}}
\newcommand{\eps}{\varepsilon}
\newcommand{\dd}{\,\mathrm{d}}
\begin{document}
\section*{UnsolvedMath ID: 30004191}
\textbf{Source problem code:} catalogue entry
\subsection{Definitions and mathematical statement}
Let \(n\geq 2\) be an integer and let \(\Omega\subset\R^n\) be a nonempty bounded open domain. No regularity assumption on \(\partial\Omega\) is supplied. For a Lebesgue-measurable set \(E\subseteq\Omega\), let \(|E|\) denote its \(n\)-dimensional Lebesgue measure. Under the conventional interpretation of perimeter in the whole ambient space, define
\[
P(E):=\sup\left\{\int_E\operatorname{div}\varphi(x)\,\dd x:\ \varphi\in C_c^1(\R^n;\R^n),\ \sup_{x\in\R^n}|\varphi(x)|\leq 1\right\}\in[0,+\infty].
\]
Thus \(E\) has finite perimeter when \(P(E)<\infty\). Define
\[
h(\Omega):=\inf\left\{\frac{P(E)}{|E|}:E\subseteq\Omega\text{ is Lebesgue-measurable},\ |E|>0\right\}.
\]
A Cheeger set of \(\Omega\) is an admissible set \(E\) satisfying \(P(E)/|E|=h(\Omega)\); measurable sets are regarded as equivalent when they differ by a null set.

\textbf{Problem.} For every such \(n\) and \(\Omega\), explicitly describe all Cheeger sets of \(\Omega\), up to null sets, and compute \(h(\Omega)\).

The supplied record does not define \(P\). The definition above uses full ambient perimeter, which includes boundary contributions along \(\partial\Omega\), rather than relative perimeter inside \(\Omega\); this interpretation requires confirmation against the original source. The record also does not specify what constitutes an explicit description or computation, or how an arbitrary domain is to be represented.
\subsection{Short English statement}
For each bounded open domain in \(\R^n\), \(n\geq2\), describe every positive-volume subset minimizing perimeter divided by volume and determine the minimum value. This draft uses full ambient perimeter; the source does not explicitly specify that convention or the meaning of an explicit computation.
\subsection{Sources}
\begin{itemize}
\item Supplied local UnsolvedMath registry record, permanent ID \(30004191\), title: \emph{Explicit Cheeger Sets and Constants of Bounded Domains}; source problem code: \texttt{catalogue entry}. No individual proposer is identified in the record.
\item Supplied original report citation: \emph{Partial Differential Equations (2020)}, Oberwolfach Reports, pp.\ 2033--2097, DOI \texttt{10.4171/owr/2019/34}. \url{https://doi.org/10.4171/owr/2019/34}. The background describes these reports as documenting mathematical workshops and their open-problem discussions. The registry records the proposed year as 2019; these supplied year labels are retained without reconciliation.
\item Supplied literature source: \emph{The Cheeger Constant of a Jordan Domain Without Necks}. \url{https://arxiv.org/abs/1704.07253}.
\item Literature assessment dated exactly \texttt{2026-08-22}: the supplied curation labels the problem partially solved and reports that explicit Cheeger sets and constants are known for important classes such as Jordan domains without necks, while no general description or computation for arbitrary domains is known. This is supplied dated curation derived from the Oberwolfach Reports record, not an independently verified literature assessment. The record requests independent verification of time-sensitive status and cited claims.
\end{itemize}
\end{document}
